\homework{1}{Wed 11 Jan 2023 10:47}{homework}

$1.3,1.5,1.8,1.12,1.17,3.3$

\begin{itemize}
	\item [1.3] Suppose that $H$ is a subgroup of $G$ of index 2 . Show that $H$ is a normal subgroup of $G$.
\begin{proof}
	From the proof of Lagrange's Theorem we know that the cosets of $H$ partitions $G$. Consider a left cost $aH \subset G$ where $a \in G$. Then either  $aH = H$ or $aH = G - H$ (set exclusion). Similarly there are only two right cosets, i.e. $Ha = H$ or $Ha = G - H$. If $aH = H$ then $a \in H$, so $Ha = H = aH$. If $aH = G - H$ then $a \not\in  H$ so $Ha \neq H$, whence $Ha = G - H = aH$. In both cases we have $aH = Ha$, so $H \triangleleft G$.
\end{proof}

	\item [1.5] Suppose that $H$ is a normal subgroup of $G$ and that $K$ is a normal subgroup of $H$. Is $K$ necessarily a normal subgroup of $G$ ?
\begin{proof}
	No. Consider the following example: We have $SL(2, \mathbb{R}) \triangleleft GL^*(2, \mathbb{R})$ as multiplicative groups. Then we consider the map
	\[
		R_3 = \left\{\left[
				\begin{array}{cc} 
			\cos\left( \frac{k}{3} \pi\right)  &- \sin\left( \frac{k}{3} \pi \right) \\
			\sin\left( \frac{k}{3} \pi \right) & \cos\left( \frac{k}{3} \pi \right) 
			\end{array}
	\right] : k \in \mathbb{Z} \right\}
	.\] 
	Then $R_3 \triangleleft SL(2, \mathbb{R})$. But it is not the case that $R_3 \triangleleft GL(2, \mathbb{R})$.
\end{proof}

	\item [1.8] Describe the elements of $\mathbb{Z}_n$ which generate $\mathbb{Z}_n$ (for any positive integer $n$ ).
\begin{proof}
	Any element $x + n\mathbb{Z} \in \mathbb{Z}_n$ such that $(x, n) = 1$ would generate $\mathbb{Z}_n$.

	If $(x, n) = 1$ then there exists $p, q \in \mathbb{Z}$ such that $p x + q n = 1$, thus $px \equiv 1 \pmod n$. Thus for any $y +  n\mathbb{Z} \in \mathbb{Z}_n$ such that $y \equiv m \pmod n$, we have $mpx \equiv m \pmod n$, hence $(mp) x + n \mathbb{Z} = y + n \mathbb{Z}$, thus $x + n\mathbb{Z}$ generates $\mathbb{Z}_n$.

	On the other hand, if $(x, n) = a \neq 1$, then $a | x$. Thus $n | \frac{n}{a} x$. This means the order of cyclic group generated by $x$ is no more than $n / a$. Hence $x$ cannot generate $\mathbb{Z}_n$ since there's not enough elements in $\left\langle x \right\rangle $.
\end{proof}

	\item [1.12] Let $\mathbb{Z}_n$ be equipped with a multiplication by defining $(a+n \mathbb{Z}) \circ$ $(b+n \mathbb{Z})=a b+n \mathbb{Z}$. Show that this is well defined, and that with this multiplication $\mathbb{Z}_n$ is a field if and only if $n$ is a prime number.
\begin{proof}


To show well-definedness, let  $a + n\mathbb{Z} = a' + n\mathbb{Z}$ and $b + n \mathbb{Z} = b' +  \mathbb{Z}$. We need to show $(ab+ n \mathbb{Z}) = (a'b' + n \mathbb{Z})$.

This is true since $ab \equiv a'b' \pmod n$, which is implied by $a \equiv a' \pmod n$ and $b \equiv b' \pmod n$.

To show $\mathbb{Z}_n$ is a field iff it is an integral domain, since finite integral domains are fields (and of course finite fields are integral domains). $\mathbb{Z}_n$ is an integral domain iff it has no zero divisor. Suppose $ab = 0$ where $a, b \in \mathbb{Z}_n$. If $n = p$ is prime then $p| a$ or $p | b$ by property of integer arithmetics, so $\mathbb{Z}_n$ has no zero divisor. If $n$ is not prime then let $n = a b$, $1 <a, b < n$. Now $ab = 0$ but $a, b \neq 0$. So $\mathbb{Z}_n$ has zero divisors.

\end{proof}

	\item [1.17] Suppose that $K$ is an infinite field and that $V$ is a vector space over $K$. Show that it is not possible to write $V=\bigcup_{i=1}^n U_i$, where $U_1, \ldots, U_n$ are proper linear subspaces of $V$.
\begin{proof}
	Suppose for the sake of contradiction that the identity is true. We may assume that $n$ is minimal, such that there is no such way to write $V$ as union of $n - 1$ proper linear subspaces.

	Now we can pick $\beta_1 \in U_1$ such that $\beta_1$ is not contained in $U_2,\ldots, U_n$. Suppose the contrary, that there does not exist  $\beta_1 \in U_1$ not contained in $U_2, \ldots, U_n$. Then we can remove $U_1$ from the union without any impact, contradicting the minimality assumption. We choose $\beta_2 \in U_2$ in the same way.

	Now consider $S = \text{span}(\beta_1, \beta_2) = K\beta_1 + K\beta_2$.
	Since $K$ is infinite, $S$ is infinite and yet $S \subset \bigcup_{i = 1} ^n U_i$. Hence by pigeonhole principle there exists distinct $s_1, s_2 \in S$ and $i \in \{1, \ldots, n\} $ such that $s_1, s_2 \in U_i$. But this implies, by algebraic closure of linear subspace, that $S \subset U_i$, and in particular, $\beta_1, \beta_2 \in U_i$, a contradiction. 
\end{proof}

	\item [3.3] Show that an integral domain with a finite number of elements is always a field.
\begin{proof}
	Let $D$ be such an integral domain, so $D$ is a commutative ring with unity $1 \in D$. For any element in $a \in D$, we have $\{a^i: i \in \mathbb{N}\} \subset   D$, which means there exists some distinct integers $i, j$ such that $a^i = a^j$. By cancellation we have  $a^{i - j} = 1$, where we assume $i - j > 0$.  Hence $a^{i - j - 1} a =  1$, so $a^{i - j - 1}$ is the multiplicative inverse of $a$. This proves that $D$ is a division ring, hence a field. 
\end{proof}
\item 1) Let $A$ be a set and define $A u t(A)=\{$ bijection: $f: A \rightarrow A\}$. Define operation
\[
A \times A \rightarrow A
\]
by the composite of bijections
\[
(f g)(a)=f(g(a)), \forall a \in A
\]
. Show that Aut(A) is a group with this operation.
\begin{proof}
	We check the group axioms.
	\begin{enumerate}
		\item (well-definedness of group operation): Composition of bijections is well-defined.
		\item (unit): Take the identity map.
		\item (inverse): Bijections are invertible.
		\item (associativity): Let $f, g, h$ be automorphisms, and we want to show $(fg) h = f(gh)$. Take an element $a \in A$. Now
			\[
				\left((fg)h\right)(a) = (fg)\left(h(a)\right) = f\left( g\left( h(a) \right)  \right) 
			.\] 
and
\[
	f(gh)(a) = f\left( (gh) (a) \right)  = f\left( g\left( h(a) \right)  \right) 
.\] 
So the associativity law holds.
	\end{enumerate}
\end{proof}

\item 2) Let $A$ be ring (recalled that $A$ is always assumed commutative and 1 exists in $A$ ) with characteristic $p$. Show that the map
\[
\varphi: x \rightarrow x^p, \forall x \in A
\]
is a ring homomorphism, which is called Frobenius.
\begin{proof}
	We only need to check that $\varphi$ preserves ring operations. Observe that
	\[
		\varphi(x)\varphi(y) = x^p y^p = (xy)^p = \varphi(xy)
	.\] 
	Now we prove that addition is preserved,
	\begin{align*}
		\varphi(x)+\varphi(y) &= x^p + y^p \\
				      &= \sum_{n=0}^{p} \binom{p}{n} x^{p - n} y^{n} \\
				      &= (x+y)^p \\
				      &= \varphi(x + y)
	.\end{align*}
	where second equality holds because $A$ has characteristic $p$.
\end{proof}

\end{itemize}

